\documentclass[11pt]{article}
    \usepackage[utf8]{inputenc}
    \usepackage[T1]{fontenc}
    \usepackage[margin=1in]{geometry}
    \usepackage{hyperref}
    \usepackage{enumitem}
    \title{Coding with code.org | (Course A) | Learning to Drag & Drop}
    \author{Piter Garcia}
    \date{March 20, 2026}
    \begin{document}
    \maketitle
    \section*{Quick Scan}
    \begin{itemize}[leftmargin=*]
    \item Grade: Kindergarten
    \item Workbook sessions: Session 20
    \item Class blocks: Day 1 12:25-1:15 Justinger, Day 2 12:25-1:15 Kesys, Day 4 12:25-1:15 Pum
    \item Reference file: coding-with-code-org-course-a-learning-to-drag-and-drop.bib
    \end{itemize}
    \section*{School Planning Goals and Inclusion}
    \begin{itemize}[leftmargin=*]
    \item What is expected: K-1.IC.2: I can explain classroom chromebook expectations | K-1.CT.4: I can decompose a task into smaller steps | K-1.CT.9: I can fix (debug) mistakes in an algorithm | K-1.CT.10: I can create a plan (algorithm) that shows steps to complete a task.
    \item What we achieved: No direct transcript match is attached yet. Treat this lesson as workbook-backed and enrich it when a matching classroom transcript is available.
    \item What needs improvement: stronger proof, cleaner assessment notes, and tighter lesson artifacts.
    \item How to improve: add transcript proof, one saved artifact, and clearer success criteria.
    \item Inclusion focus: use visual modeling, chunked directions, predictable routines, and flexible participation roles.
    \end{itemize}
    \section*{Official References}
    \begin{itemize}[leftmargin=*]
    \item \url{https://csteachers.org/k12standards/interactive/}
\item \url{https://code.org/curriculum/elementary-school}

    \end{itemize}
    \end{document}
